%!TEX program = lualatex
\documentclass[ngerman,headheight=70pt]{scrartcl}
\usepackage[ngerman]{babel}
\usepackage{blindtext}
\usepackage[utf8]{luainputenc}
\usepackage{scrlayer-scrpage}
\usepackage[T1]{fontenc}
\usepackage{fontspec}
\usepackage{multicol}
\usepackage{enumitem}
\usepackage{epigraph}
\usepackage[german=quotes]{csquotes}
\usepackage[bottom=2cm,footskip=8mm]{geometry}
\setmainfont{Source Sans Pro}
\MakeOuterQuote{"}

\setlength{\parindent}{0pt}
\setlength{\parskip}{6pt}

\pagestyle{plain.scrheadings}
\thispagestyle{scrheadings}
\lohead{\vspace*{-1.5cm}\includegraphics[scale=0.4]{../images/stupa_logo_rot.pdf}}
\rohead{\\Präsidium des\\Studierendenparlaments}

\lofoot[Universität Hamburg · Präsidium des Studierendenparlaments\\\hrule Von-Melle-Park 5 · D-20146 Hamburg]{Universität Hamburg · Präsidium des Studierendenparlaments\\\hrule Von-Melle-Park 5 · D-20146 Hamburg}
\cofoot[]{}
\cohead[\thepage]{}

\renewcommand{\thesubsection}{TOP \arabic{subsection}}
\renewcommand{\thesubsubsection}{\arabic{subsubsection}.}
\setcounter{subsection}{-1}

\begin{document}
    UHH · StuPa-Präsidium · Von-Melle-Park 5 · D-20146 Hamburg

    \section*{Beschluss des Studierendenparlamentes vom 30. Juni 2016}
    \subsection*{Weiter initiativ mit der Strategie "Uni für alle statt Exzellenz"!}

    \blockquote[Jochen Hörisch, "Privat ein Laster, öffentlich eine Tugend",
    FAZ vom 29. Mai 2016]{"Der Pressesprecher oder Präsident einer Universität
    verkündet stolz, dass seine Universität im Ranking in den letzten Jahren von
    Platz 37 auf Platz 23 aufgestiegen sei, wird aber sprachlos, wenn er berichten
    soll, welchen großartigen Einsichten und Forschungsergebnissen dieser Aufstieg
    denn zu verdanken sei. Beredt wird er wieder, wenn er kundtut, man habe soundso
    viele Graduiertenkollegs gegründet und einen Sonderforschungsbereich eingerichtet,
    aus der Exzellenzinitiative größere Summen erhalten etc. pp. Im Übrigen gebe
    es noch viel zu tun, die Universität werde sich noch um weitere Fördergelder
    bewerben, man sei zuversichtlich, diese zu erhalten, denn die Universität sei
    gut aufgestellt."}

    Die Exzellenzinitiative heißt jetzt Exzellenzstrategie, und bleibt weiterhin
    eine große Dummheit. Nach wie vor schadet sie erheblich Forschung und Lehre,
    quält die Hochschulangehörigen und konterkariert eine gesellschaftlich
    verantwortungsvolle Wissenschaft im Dienste der der großen Mehrheit der
    Menschen.

    Entgegen nachdrücklicher und immer weitere Kreise ziehende Kritik und Proteste
    aus den Hochschulen haben sich die Ministerpräsidenten und die Bundeskanzlerin
    am 16. Juni auf die Fortsetzung des "Exzellenz"-irgendwas geeinigt -- mit ein
    wenig mehr Geld und verbunden mit dem Gefasel von "besseren Aufstiegschancen".
    Betont wird ansonsten das pur ideologische Gebot eines von "strengem
    wissenschaftlichen Wettbewerb geleiteten Verfahrens". Es ist wenig überzeugend,
    wenn die Ministerpräsidentenkonferenz ständig hervorheben muss, wie wissenschaftlich
    ihre Wissenschaftspolitik ist.

    Es ist also noch nicht gelungen, eine nächste Runde der Exzellenzinitiative
    jetzt schon zu verhindern zu Gunsten einer bedarfsdeckenden Grundfinanzierung
    aller Hochschulen. Dennoch haben die bundesweiten Proteste die öffentliche
    Diskussion zugunsten einer Durchsetzung humanistischer Bildung und Wissenschaft
    in Überwindung der üblen neoliberalen Deformierungen deutlich vorangebracht.
    Die demonstrative Euphorie für das Elite-Ideologem ist einem zunehmenden
    Realismus gewichen, die Wahrheit über die Exzellenzinitiative wird immer
    häufiger offen ausgesprochen: Mit inhaltlicher Qualität hat sie nichts zu tun;
    die Konkurrenzförmigkeit schadet einer kooperativen Entwicklung von Forschung
    und Lehre mit dem Inhalt einer solidarischen gesellschaftlichen Entwicklung;
    das Antragswesen verschleißt nutzlos enorme Kräfte, die einer sinnvollen
    wissenschaftlichen Produktivität geraubt werden; die Antragsprosa ist eine
    einzige Heuchelei und Schönfärberei, die den argumentativen Austausch über
    ein vernünftige Weiterentwicklung der Wissenschaft verhindert.

    Der Kaiser steht schon lange ohne Textilien da, eine neue Qualität der
    Entzauberung des Despoten ist erreicht, weil immer mehr dem Kind beipflichten,
    dass der Despot wirklich nackt ist -- nun muß er noch entthront werden.

    Das Studierendenparlament und der AStA werden daher die Kampagne "Uni für
    Alle statt Exzellenzinitiative" weiterführen und dafür die fortgesetzte und
    erweiterte Zusammenarbeit mit Fachschaftsräten und Mitgliedern aller
    Hochschulgruppen suchen.

    \begin{itemize}
        \item Die Unterschriftenkampagne wird textlich aktualisiert und fortgeführt
        \item In allen Gremien wird die Debatte um Uni für alle statt Exzellenz
        weitergeführt, darauf gerichtet, dass keinerlei strukturelle und
        finanzielle Unterwerfung unter die Exzellenzinitiative stattfindet, und
        mit der Erweiterung auf eine generelle Kritik an dem Konkurrenzprinzip
        in der Wissenschaft, speziell an der Finanzierung durch Drittmittel.
        \item Als positiver Gegenpol wird die Studienreform weiter voran getrieben
        mit verstärktem Augenmerk auf die kritische Verantwortung der Wissenschaft
        für Frieden, soziale Gerechtigkeit, ökologische Nachhaltigkeit, Demokratie
        und Flüchtlingssolidarität.
        \item Das Wintersemester wird mit Aktivitäten (u.a. Vollversammlung, Demo)
        zur Beendigung der Schuldenbremse für eine bedarfsdeckende Finanzierung
        von Hochschule und Bildung, sowie anderen sozialstaatlichen Einrichtungen
        eröffnet. Dafür wird die Zusammenarbeit mit BündnispartnerInnen auch
        außerhalb der Hochschulen gesucht.
    \end{itemize}
\end{document}
